% ===== §15 =====
\section{Envoi: A Promise That Is Not Optimized}
\label{p25:sec:envoi}

\noindent
The argument has run its course, from the anatomy of the optimizing frame, through the four impossibilities that show it ill-posed in intimate relation, through the criticism of its imposition, to the recovery of a rationality that responds rather than maximizes. It remains to close as the series closes its papers, not with an argument but with an image, and to return, as the paper began, to the plainest form of the thing it has been about.

Consider a promise. Two people bind themselves to one another for a life, and in doing so they perform an act the optimizing frame cannot understand. The frame would have them compute: to weigh the expected value of the bond against its alternatives, to keep the option of revision open should a better prospect appear, to treat the vow as a provisional maximum held only until the numbers change. A promise made in that spirit is not a promise but a standing intention to defect when it pays, and everyone knows the difference, though the frame cannot state it. For a promise is precisely the refusal to keep computing. It is the act by which one closes the ledger, forecloses the revision, and binds oneself past the point where the calculation could reach, saying not ``I will stay while you remain the optimal choice'' but ``I will stay,'' full stop, whatever the sums come to say. The whole worth of a promise lies in its being unconditional in just the way an optimization can never be, and the series has argued elsewhere that this unconditionality, far from a weakness, is the very source of the trust a promise can generate, since it is the mark that the one who promises has placed the bond beyond the reach of the calculation that could betray it \parencite{paperxiii}.

This is why a promise is a promise only when it is not optimized, and why the attempt to make it so destroys it in the same motion by which the pricing destroyed the friendship and the maximizing settled the circulation. To ask, of the one you have bound yourself to, whether she is worth it, to hold the vow open against the day a better yield appears, is not to be a more rational lover but to have stopped promising and started calculating, and to have lost, in the calculating, the very good the promise was. The good of the promise is available only on the far side of the refusal to compute its worth; it is a good one can have only by declining to ask the question the optimizing frame insists is the rational one. And so the deepest thing the paper has to say is said not by its arguments but by this plainest of acts, in which one human being says to another that here the question of worth will not be asked, not because the answer might be unwelcome but because the asking would already be the betrayal. There are goods a love does not maximize, and it is not for want of a better measure. It is because to have reached for the measure at all would have been to have loved something other than the one before you, who is not a quantity of value that a greater quantity could replace, but this one, irreplaceable, whom one does not weigh. \zh{未曾计算，故其爱得全}---never weighed, and so the love is whole; which is to say, a promise, and not a sum.
